\section{Comparison with contrary flatness claims}\label{sec:concentration}

Corollary~\ref{cor:finite-flatness} conflicts with several stated
nonflatness results. We compare the precise versions listed in the
bibliography and examine the proposed obstructions. This discussion is
independent of the construction proving Theorem~\ref{thm:main}.

An earlier preprint of el Abdalaoui asserted nonexistence of
square-$L^2$-flat real Littlewood sequences
\cite[Theorem~3.3]{ElAbdalaoui2017}. Its definition fixes both endpoint
signs to $+1$ \cite[Equation~(2.1)]{ElAbdalaoui2017}. Changing at most
two endpoint signs of our length-$N$ polynomials costs at most
$4/\sqrt N$ in normalized uniform norm. Their normalized maxima still
tend to one, so the proof of Corollary~\ref{cor:finite-flatness} gives
square-$L^2$ flatness with this convention as well.

\subsection{The concentration constant and a neighborhood of zero}
Theorem~1 of \cite{ElAbdalaoui2025} claims that $L^2$-normalized real-sign
polynomial sequences cannot be $L^{2p}$-flat for integers $p>1$. This conflicts with
Corollary~\ref{cor:finite-flatness}, even though our maximum estimate
supplies no uniform lower bound. The final inference in the cited proof
uses a concentration constant with a different quantifier order from
the one needed there.

To see the distinction, let $\mathcal I$ be the nonzero analytic
idempotent trigonometric polynomials
\[
 I(t)=\sum_{k\in E_0}\e(kt),
 \qquad \varnothing\ne E_0\subset\Z_{\ge0}\text{ finite}.
\]
For $\alpha>0$ and a measurable set $E\subset\T$, define
\[
 B_\alpha(E)=\sup_{I\in\mathcal I}
      \frac{\int_E|I(t)|^\alpha\,dt}{\int_{\T}|I(t)|^\alpha\,dt}.
\]
The concentration level of Bonami--R\'ev\'esz is
\cite[Definitions 1--2]{BonamiRevesz2008}
\begin{equation}\label{eq:concentration-infimum}
 c_\alpha=\inf_{\substack{E\text{ symmetric, open}\\E\ne\varnothing}}
                B_\alpha(E).
\end{equation}
Indeed every number strictly smaller than the infimum is a concentration
guarantee for each such $E$, and every uniform guarantee is no larger
than the infimum. Thus an upper bound on $c_\alpha$ does not give that
upper bound on $B_\alpha(E)$ for a prescribed $E$.

For completeness, the particular sets used in the contrary inference
have optimal concentration equal to one. Suppose $E$ contains the
neighborhood $(-d,d)$ of zero, where $0<d<1/2$, and put
$D_q(t)=\sum_{k=0}^{q-1}\e(kt)$. The geometric-series formula gives
\[
 |D_q(t)|\le\frac1{\sin(\pi d)}\quad(t\notin E).
\]
For $|t|\le1/(2q)$, the same formula and
$\sin x\ge2x/\pi$ on $[0,\pi/2]$ give
$|D_q(t)|\ge2q/\pi$. Consequently, for $\alpha>1$,
\begin{align*}
 \int_{\T\setminus E}|D_q|^\alpha&\le\sin(\pi d)^{-\alpha},\\
 \int_{\T}|D_q|^\alpha&\ge(2/\pi)^\alpha q^{\alpha-1}.
\end{align*}
Their ratio tends to zero, proving $B_\alpha(E)=1$. Bonami--R\'ev\'esz
also explicitly distinguish full concentration at zero from the uniform
concentration level \cite[Theorem 7 and Proposition 9]{BonamiRevesz2008}.

Under its assumed flatness, Equation~(27) of \cite{ElAbdalaoui2025}
gives concentration tending to
one in each fixed neighborhood of zero. This is compatible with
\eqref{eq:concentration-infimum}; it does not contradict the cited upper
bound $c_{2p}\le1/2$. Multiplicative normalization cancels in the ratio.
Even restricting supports to $\{0,\ldots,q-1\}$ with density tending to
$1/2$ does not give the required cap: use $D_{\lfloor q/2\rfloor}$.
This explains the failure of that specific concluding inference.

\subsection{A later weighted assertion}
A separate preprint \cite[Theorem 1]{ElAbdalaoui2025Generalization}
asserts a nonflatness criterion for fixed real coefficients $a_j$ and
unimodular multipliers. Its hypothesis is
\begin{equation}\label{eq:weighted-hypothesis}
 \sum_{j=1}^n a_j^2\le\frac K{n^2}\sum_{j=1}^n j^2a_j^2,
\end{equation}
and its polynomial is normalized by $(\sum_{j=1}^n a_j^2)^{1/2}$.
Without any restriction preventing a dominant coefficient, this
criterion is false.

Take $a_j=2^{j^2}$ and all unimodular multipliers equal to one. Let
$S_n=\sum_{j=1}^n a_j^2$. For $n\ge2$,
\[
 \frac{\sum_{j<n}a_j^2}{a_n^2}
 \le(n-1)2^{-4n+2}\le1.
\]
It follows, also at $n=1$, that
\[
 S_n\le2a_n^2\le\frac2{n^2}\sum_{j=1}^n j^2a_j^2.
\]
Thus \eqref{eq:weighted-hypothesis} holds with $K=2$. On the other hand,
\[
 \frac{\sum_{j<n}a_j}{a_n}\le(n-1)2^{-2n+1}\longrightarrow0,
 \qquad \frac{S_n}{a_n^2}\longrightarrow1.
\]
Uniformly for $|z|=1$, therefore,
\[
 \left|\frac{\sum_{j=1}^n a_jz^j}{\sqrt{S_n}}-z^n\right|
 \le\frac{\sum_{j<n}a_j}{\sqrt{S_n}}
       +\left|\frac{a_n}{\sqrt{S_n}}-1\right|\longrightarrow0.
\]
The normalized moduli converge uniformly to one. This example concerns
weighted coefficients, not signs.

The proof of Corollary~3 in
\cite{ElAbdalaoui2025Generalization} displays a signed-modulus identity
for values at $|z|=|z'|=1$.
Even reading its repeated second argument as the intended $z'$, the
asserted form
\[
 |P(z)-P(z')|=\sigma|P(z)|+\tau|P(z')|,
 \qquad \sigma,\tau\in\{-1,1\},
\]
is false. For $P(z)=1+z$, $z=1$ and $z'=i$, the left side is $\sqrt2$,
whereas no value $\pm2\pm\sqrt2$ equals $\sqrt2$.
These observations identify defects in the proposed obstructions;
the real-sign assertion of Theorem~\ref{thm:main} rests on the
construction and rounding proof above.
